% Annexe du bloc « fiscalité immobilière » de la note de 8 pages (note8p/CADRE.md). Générée par
% note8p/calculs/foncier_annexe.py à partir de note8p/sorties/foncier.json (note8p/calculs/foncier_calculs.py, qui importe
% foncier/simulation.py). Compte rendu : note8p/sorties/foncier.md. Montants en Md€ de 2025, assiettes 2025 non indexées.

\subsection{Données, champ et calibrage}

Le modèle communal est celui de la note fusionnée (\texttt{foncier/simulation.py}, inchangé) : 30,5~millions de logements
des particuliers de métropole hors Alsace-Moselle ; valeur = prix médian au m² des ventes 2023-2024 de la commune × surface
médiane ; part du sol mesurée par les terrains à bâtir (maisons) ou par résidu (appartements), calée sur les comptes de
patrimoine des ménages (46,2~\% fin 2025) ; taxe foncière de 2025 (REI) répartie au prorata de la valeur dans chaque
commune. Les agrégats nationaux viennent des comptes de patrimoine de l'Insee (tableau 8.200, fin 2025) : biens privés
(sociétés non financières, sociétés financières, ménages, institutions sans but lucratif), hors terres agricoles et biens
publics. Sol : 6\,656~Md€, dont 565 de terrains non bâtis ; bâti (logements et bâtiments non résidentiels)
: 6\,560~Md€.

\textbf{Droits de mutation.} Taux permanent de 5,807~\% : 4,50~\% pour le département (art.~1594~D CGI), 1,20~\% pour la
commune (art.~1584) et 2,37~\% de la part départementale de frais d'assiette (art.~1647~V). La majoration à 5~\% (loi
n°~2025-127, art.~116), dont les primo-accédants sont exclus, cesse le 31~mars 2028 ; la référence de la réforme est donc
le taux permanent. France entière, au taux permanent et sur les ventes de 2025, les droits valent 13,1~Md€
sur les logements et 3,0~Md€ sur les autres biens (différence OFGL moins Compte du logement, ramenée au
taux permanent), soit 16,1~Md€ : la division par deux coûte \textbf{8,07~Md€}, contre
7,58 dans l'ancien calibrage, limité au champ du modèle. Le coût suit le cycle :
7,2~Md€ aux ventes moyennes de 2023-2024, 10,0~Md€ à celles de 2022. Un point de taux vaut
2,78~Md€.

\subsection{Architecture et formules}

Chaque bien $i$ a une base $b_i = S_i + B_i/4$ ($S$ : valeur du sol, bâti ou non ; $B$ : valeur du bâti). Il paie
$(t_c + t_n)\,b_i$ : le sol est taxé au taux $t$ et le bâti au taux $t/4$. Lors de la bascule, le taux communal
$t_c = \sum_{i \in c} \mathrm{TF}_i^0 / \sum_{i \in c} b_i$ garde la recette de chaque commune ; la commune le vote
ensuite librement, le rapport de 4 étant fixé par la loi. Le taux national uniforme $t_n$ lève l'enveloppe $E$ :
$t_n = E / \sum_i b_i$, la part des logements des ménages (64,5~\% de la base) étant répartie
entre les logements du modèle au prorata de $b_i$.

\textbf{Capitalisation.} La valeur du sol est la rente nette capitalisée : $S = \rho/(r + \tau_S)$, où $\tau_S$ est
l'impôt sur le sol rapporté à sa valeur. Aujourd'hui, $\tau_S^0 = 0{,}358~\%$ (taxe / valeur des logements, appliquée à
tous les biens bâtis) et 0,08~\% pour les terrains non bâtis (hypothèse). Après la réforme, $\tau_S^1 = t_c + t_n$ ; le
taux communal monte pour garder la recette sur un sol moins cher. Le point fixe
$S^1 = S^0 (r+\tau_S^0)/(r+\tau_S^1)$, $t_n = E/(S^1 + B/4)$ est résolu numériquement avec $r = 3~\%$ (variante 4~\%), le
bâti restant à son coût. La perte de valeur des biens est la valeur actuelle de la hausse, $\Delta \mathrm{TF}_i / r$
(baisse des droits gardée par l'acheteur, cas central ; variante symétrique dans le tableau~\ref{tab:fon8_env}).

\begin{table}[!htbp]\centering\scriptsize
\caption{Enveloppes de recette foncière nouvelle}\label{tab:fon8_env}
\begin{tabular}{@{}lrrrrrrrrr@{}}
\toprule
Enveloppe & \multicolumn{3}{c}{Taux national sur le sol (\%)} & Baisse du & Recette au & Payé par & \multicolumn{2}{c}{Perte de valeur (Md€)} & Net \\
\cmidrule(lr){2-4}\cmidrule(lr){8-9}
(Md€) & statique & $r$ = 3 \% & $r$ = 4 \% & sol (\%) & taux statique & logements & tous & logements & 2040 \\
\midrule
8,1 (substitution) & 0,097 & 0,106 & 0,103 & 9,8 & 7,4 & 5,2 & 269 & 173 & $-$0,8 \\
10 & 0,121 & 0,132 & 0,129 & 10,6 & 9,2 & 6,4 & 333 & 215 & +1,0 \\
15 (recommandée) & 0,181 & 0,201 & 0,196 & 12,4 & 13,5 & 9,7 & 500 & 322 & +5,5 \\
20 & 0,241 & 0,272 & 0,264 & 14,3 & 17,7 & 12,9 & 667 & 430 & +10,0 \\
\bottomrule
\end{tabular}
\par\vspace{2pt}{\raggedright\scriptsize Taux du bâti : un quart du taux du sol. « Recette au taux statique » : ce que
rapporterait le taux calculé sans capitalisation, une fois les valeurs ajustées. Perte de valeur à $r$ = 3~\% (4~\% : un
quart de moins). Si la baisse des droits se capitalise aussi, la perte nette de tous les propriétaires vaut
231~Md€ pour 15~Md€. Sans les terrains non bâtis, tous les taux montent de 7~\%.\par}
\end{table}

\subsection{Résultats pour les logements}

\begin{table}[!htbp]\centering\scriptsize
\caption{Taxe foncière des logements par dixième de prix au m² de la commune}\label{tab:fon8_dec}
\begin{tabular}{@{}rrrrrrrrrr@{}}
\toprule
 & Prix & Taxe/valeur & \multicolumn{4}{c}{Variation de taxe (€/an ; \% pour 15)} & Taxe $-$ droits & Taxe/valeur & Valeur \\
\cmidrule(lr){4-7}
Dixième & (€/m²) & 2025 (\%) & 10 & 15 & 15 (\%) & 20 & 15 (€/an) & 15 (\%) & 15 (\%) \\
\midrule
1 & 1\,035 & 0,54 & +36 & +60 & +13 & +84 & $-$21 & 0,61 & $-$2,4 \\
2 & 1\,404 & 0,49 & +61 & +97 & +17 & +132 & $-$9 & 0,57 & $-$2,7 \\
3 & 1\,693 & 0,49 & +79 & +124 & +18 & +169 & $-$6 & 0,58 & $-$2,9 \\
4 & 2\,000 & 0,46 & +125 & +184 & +24 & +243 & +38 & 0,57 & $-$3,7 \\
5 & 2\,356 & 0,43 & +161 & +235 & +28 & +310 & +70 & 0,56 & $-$4,1 \\
6 & 2\,737 & 0,41 & +185 & +274 & +31 & +363 & +87 & 0,54 & $-$4,3 \\
7 & 3\,176 & 0,39 & +233 & +340 & +37 & +448 & +131 & 0,54 & $-$4,8 \\
8 & 3\,749 & 0,36 & +278 & +413 & +42 & +548 & +182 & 0,51 & $-$5,0 \\
9 & 4\,746 & 0,31 & +331 & +496 & +49 & +662 & +202 & 0,46 & $-$5,0 \\
10 & 6\,261 & 0,18 & +513 & +780 & +93 & +1\,047 & +310 & 0,35 & $-$5,6 \\
\bottomrule
\end{tabular}
\par\vspace{2pt}{\raggedright\scriptsize Enveloppes de 10, 15 et 20~Md€ en régime permanent ; « taxe $-$ droits » :
hausse de taxe moins la baisse des droits rapportée à l'ensemble des logements de la commune ; valeur : perte de valeur
capitalisée ($r$ = 3~\%). Avec 15~Md€ : hausse médiane +237~€, 9\textsuperscript{e}
décile +655~€, 99\textsuperscript{e} centile +975~€ ;
40~\% des logements au-delà de +300~€, 12~\% au-delà de +600~€ ;
baisse pour 5~\% ; Paris +156~\% ; maisons
+263~€, appartements +351~€. Impôt sur le bâti des logements :
14,3~Md€ aujourd'hui, 7,3 avec la seule part communale à la nouvelle base, 9,4 avec 15~Md€.\par}
\end{table}

\subsection{Trajectoire 2027-2040}

Calendrier (choix) : loi en 2027, valeurs de référence au 1\textsuperscript{er}~janvier 2028, premier avis à l'automne
2029 ; part nationale et passage de la part communale à la nouvelle base par tiers en 2029, 2030 et 2031 ; taux des droits
abaissé au 1\textsuperscript{er}~janvier des mêmes années. Une clause suspend chaque palier de baisse tant que le palier
de taxe n'est pas mis en recouvrement. Hypothèses : 3~\% des ventes annuelles décalées de décembre à
janvier par point de baisse annoncée (les +26~\% de ventes du mois précédant la hausse de 0,7~point de 2014 représentent
environ 2~\% des ventes de l'année, soit 3~\% par point) ; report ouvert aux propriétaires occupants modestes, qui doivent
25~\% des hausses des occupants, utilisé par la moitié d'entre eux, avec un intérêt de 3~\% et un
remboursement de 1/15 du stock par an ; déductibilité l'année suivante au taux effectif de 20~\% pour
la part des sociétés soumises à l'IS (hors HLM, 18,7~\% de la base) et de
25~\% pour la hausse des bailleurs. Aucune reprise durable des ventes n'est comptée : à
6~\% de ventes par point, le coût des droits tomberait de 8,1 à 6,6~Md€.

\begin{table}[!htbp]\centering\scriptsize
\caption{Trajectoire budgétaire (Md€ ; + = recette publique)}\label{tab:fon8_traj}
\begin{tabular}{@{}lrrrrrrrrr@{}}
\toprule
 & 2027 & 2028 & 2029 & 2030 & 2031 & 2032 & 2033 & 2035 & 2040 \\
\midrule
Taxe foncière nouvelle & 0 & 0 & +5,0 & +10,0 & +15,0 & +15,0 & +15,0 & +15,0 & +15,0 \\
Droits de mutation & 0 & 0 & $-$2,7 & $-$5,4 & $-$8,1 & $-$8,1 & $-$8,1 & $-$8,1 & $-$8,1 \\
Achats différés & 0 & $-$0,5 & 0 & 0 & +0,2 & 0 & 0 & 0 & 0 \\
Report de paiement & 0 & 0 & $-$0,2 & $-$0,4 & $-$0,6 & $-$0,5 & $-$0,5 & $-$0,4 & $-$0,3 \\
Déductibilité & 0 & 0 & 0 & $-$0,4 & $-$0,8 & $-$1,2 & $-$1,2 & $-$1,2 & $-$1,2 \\
\textbf{Net, 15 Md€} & 0 & $-$0,5 & +2,1 & +3,8 & +5,8 & +5,2 & +5,2 & +5,3 & +5,5 \\
Net, 10 Md€ & 0 & $-$0,5 & +0,5 & +0,8 & +1,2 & +0,8 & +0,8 & +0,8 & +1,0 \\
Net, 20 Md€ & 0 & $-$0,5 & +3,7 & +6,9 & +10,3 & +9,6 & +9,7 & +9,8 & +10,0 \\
Net, substitution & 0 & $-$0,5 & $-$0,1 & $-$0,4 & $-$0,5 & $-$0,9 & $-$0,9 & $-$0,9 & $-$0,8 \\
Net, 15 Md€, baisse unique 2030 & 0 & 0 & +3,4 & +1,8 & +5,6 & +5,2 & +5,2 & +5,3 & +5,5 \\
Net, 15 Md€, un an de retard & 0 & 0 & $-$0,5 & +2,1 & +3,8 & +5,8 & +5,2 & +5,3 & +5,4 \\
Stock de report (15 Md€) & 0 & 0 & 0,2 & 0,6 & 1,2 & 1,8 & 2,3 & 3,4 & 5,6 \\
Taux des droits (\%) & 5,81 & 5,81 & 4,84 & 3,87 & 2,90 & 2,90 & 2,90 & 2,90 & 2,90 \\
\bottomrule
\end{tabular}
\par\vspace{2pt}{\raggedright\scriptsize Années 2034, 2036-2039 : voir \texttt{note8p/sorties/foncier.json}, table
\texttt{annuel} (15~Md€) et \texttt{annuel\_par\_scenario}. Le report est un prêt remboursable, compté ici en besoin de
financement net.\par}
\end{table}

\subsection{Cas-types}

\begin{table}[!htbp]\centering\scriptsize
\caption{Achats et taxe foncière (prix médian DVF au m² de la commune × surface)}\label{tab:fon8_cas}
\setlength{\tabcolsep}{3pt}
\begin{tabular}{@{}lrrrrrrrrrr@{}}
\toprule
 & Prix & Taux & Droits & Droits & Droits & Taxe & \multicolumn{2}{c}{Hausse (€/an)} & Bilan & Perte du \\
\cmidrule(lr){8-9}
 & (€) & 2026 (\%) & 2026 & permanents & après & actuelle & 15 & 10 & 7 ans & vendeur \\
\midrule
Orléans, app. 65 m² & 170\,000 & 5,81 & 9\,871 & 9\,871 & 4\,936 & 1\,102 & +298 & +223 & +3\,079 & $-$9\,934 \\
Lyon, app. 80 m² & 372\,000 & 6,32 & 23\,505 & 21\,601 & 10\,800 & 1\,242 & +665 & +450 & +6\,660 & $-$22\,151 \\
Massy, app. 65 m² & 269\,000 & 5,81 & 15\,620 & 15\,620 & 7\,810 & 925 & +475 & +326 & +4\,849 & $-$15\,844 \\
Ste-Geneviève, maison 100 m² & 365\,000 & 6,32 & 23\,063 & 21\,194 & 10\,597 & 1\,357 & +562 & +379 & +7\,098 & $-$18\,722 \\
\bottomrule
\end{tabular}
\par\vspace{2pt}{\raggedright\scriptsize Orléans et Massy : primo-accédants (hors majoration) ; Lyon et
Sainte-Geneviève-des-Bois : autres acheteurs, à 6,32~\% jusqu'au 31~mars 2028 (tableau DGFiP au 1\textsuperscript{er}~juin
2026). Bilan à 7 ans : économie de droits moins valeur actuelle de sept ans de hausse, prix inchangés. Perte du vendeur :
hausse capitalisée à 3~\%, qui revient à l'acheteur sous forme d'un prix plus bas si la taxe se capitalise. Mobilité : une
famille qui quitte Orléans pour un poste en Essonne paie 21\,194~€ de droits au taux permanent
(23\,063 en 2026), 10\,597~€ après réforme, soit l'équivalent de 294~€
puis 147~€ d'augmentation nette par mois pendant six ans.\par}
\end{table}

\subsection{Gains, sensibilités et variantes écartées}

\textbf{Perte sèche} (indicateur d'efficacité, triangle de Harberger, méthode de l'annexe de la note fusionnée). Ventes de
logements : $-$0,86, $-$1,71 et $-$2,29~Md€
pour 3, 6 et 8~\% de ventes en moins par point (soit +9, +18 et
+25~\% de ventes durables) ; zéro si la réaction n'est que temporaire. Autres immeubles :
$-$0,40~Md€ (transposition). Bâti des logements (élasticité 1, coût d'usage 4,5~\%) :
$-$0,41, $-$0,37 et $-$0,33~Md€ pour 10, 15 et 20~Md€. Le sol ne crée
pas de perte sèche. Schéma g6 : un terrain de 200\,000~€ et un immeuble de 600\,000~€ ; au taux unique de
0,358~\%, construire ajoute 2\,149~€ d'impôt par an ; avec
0,651~\% sur le sol et 0,163~\% sur le bâti (même recette), 977~€,
et le terrain nu paie 1\,302~€ au lieu de 716.

\textbf{Variantes écartées.} \emph{Substitution pure} (8,1~Md€) : le solde est négatif ($-$0,8~Md€ par an) à cause de la
déductibilité, et rien ne finance l'IS. \emph{Taux national sur le seul sol} : il évite les 20~\% de la part nationale
qui portent sur le bâti, mais le rapport sol/bâti varierait alors entre communes et il faudrait deux bases dans l'avis.
\emph{Baisse des droits en une fois} (2030) : aucune année négative dans le calcul, mais une incitation à différer les
achats pendant plus de deux ans que l'hypothèse d'un mois de décalage ne capte pas. \emph{Taux départemental} : de 0,13
à 0,79~\% du sol selon le département (ancien calcul). \emph{Suppression des droits} : 16,2~Md€ ; orientation de long
terme seulement. \emph{Baisse immédiate avant la taxe} (ancienne trajectoire) : 8,1~Md€ de dette par année d'avance.

\textbf{Expériences de référence.} \emph{Bade-Wurtemberg} (Landesgrundsteuergesetz du 4/11/2020, modifié en 2023 et
le 18/11/2025) : valeur = surface × \emph{Bodenrichtwert} de la zone, fixé par les comités d'experts (§~38 al.~1-2) ;
preuve d'une valeur inférieure de plus de 30~\% par expertise (§~38 al.~4) ; taux légal de 1,30~pour mille, réduit de 30~\% pour
les terrains majoritairement résidentiels, de 25~\% pour le logement social, de 10~\% pour les monuments (§~40) ;
\emph{Hebesatz} unique par commune (§~50 al.~4), taux distinct possible sur les terrains constructibles non bâtis
(§~50a) ; valeurs au 1/1/2022, impôt dû en 2025 ; modèle validé par le Bundesfinanzhof (II~R~26/24 et II~R~27/24,
communiqué du 20/05/2026), recours constitutionnels pendants (source associative). \emph{Pennsylvanie} : rapports
municipaux sol/bâti de 1,4 à 17,5 dans 18 villes en 2000, dont Pittsburgh (1913-2001), Scranton (1913), Harrisburg (1975),
Allentown (1997) (Banzhaf et Lavery, 2010, tableau~2) ; comté et district scolaire taxent au taux unique. Oates et Schwab
(1997) : permis réels +70~\% à Pittsburgh entre 1960-79 et 1980-89 contre $-$14~\% en moyenne dans 15 villes, attribués
surtout à la pénurie de bureaux (comparaison avant-après, un seul traité). Plassmann et Tideman (2000) : effet positif
sur les permis dans 15 villes à deux taux contre 204 témoins, sensible à la spécification (résumé et thèse lus). Banzhaf
et Lavery (2010) : différences premières avec effets fixes par secteur de recensement ; +5 à 6 points de croissance
décennale du nombre de pièces, surtout par plus de logements ; environ +3 points (pièces) et +2 points (logements) pour un
rapport de 4 (version de travail de 2008). Pittsburgh a abandonné le double taux en 2001 après une réévaluation qui avait
relevé la valeur du sol de 81~\% (Hughes, 2006).

\textbf{Limites.} Taux uniforme de la valeur dans chaque commune (la dispersion réelle multiplie les perdants) ; valeurs
de l'Insee sans géographie pour les sociétés ; contenu des « autres terrains » ; droits sur les autres biens par
différence de deux sources ; paramètres du report et de la déductibilité supposés ; capitalisation complète supposée
(contestée par une étude danoise non publiée) ; premier avis en 2029 ambitieux (le Bade-Wurtemberg a mis plus de quatre
ans en partant de valeurs de zone existantes).
